% ECONOMICS_PROBLEM: {"id": "C78-352", "title": "Minimum reachable-state count for doubly prime stable-marriage instances", "jel_code": "C78", "source_id": "OP-0204"}
% FIRST_POSTED: 2026-10-10
% ATTEMPT_STATUS: solved
% ENTRY_KIND: research_attempt
\documentclass[11pt]{article}
\usepackage[a4paper,margin=27mm]{geometry}
\usepackage{amsmath,amssymb,amsthm,mathtools}
\usepackage[T1]{fontenc}
\usepackage{lmodern}
\usepackage[colorlinks=true,allcolors=blue]{hyperref}
\usepackage{microtype}
\newtheorem{theorem}{Theorem}
\newtheorem{lemma}[theorem]{Lemma}
\newtheorem{corollary}[theorem]{Corollary}
\newcommand{\Reach}{\operatorname{Reach}}
\newcommand{\Id}{\mathcal J}
\title{The minimum number of deferred-acceptance states in a doubly prime instance}
\author{Ji Ho Bae}
\date{10 October 2026}

\begin{document}
\maketitle
\begin{abstract}
For every $n\ge2$, the minimum number of reachable serial
deferred-acceptance states among doubly prime strict complete
stable-marriage instances of order $n$ is $3\cdot2^{n-1}-1$.
The lower bound follows by selecting one rejection witness for each
nonfirst proposal. The resulting connected directed acyclic graph has
two incoming neighbors at each nonsource, and every one of its order
ideals is a reachable state. A counting lemma for such graphs supplies
the bound. A rejection-chain family attains it with $2n-1$ proposals.
This proves the conjecture in Ishida's Section~7(A).
\end{abstract}

\section{The question and its conventions}

An instance has $n$ proposers $P$, $n$ receivers $R$, and strict complete
preference lists. A serial deferred-acceptance state is $(h,q)$, where
$q(p)$ counts proposals made by $p$ and $h(r)$ is the proposer held by
$r$, or $\bot$. Initially all counts are zero and all receivers are empty.
A legal step selects an unheld proposer with an untried receiver and
makes its next proposal. The receiver keeps its preferred of the new
proposer and its previous holder. Write $\Reach(I)$ for all states
reached by all finite legal schedules, including the initial and terminal
states.

An \emph{input-independent partition} pairs partitions of $P$ and $R$
into nonempty equally sized blocks, with every agent ranking all members
of its opposite-side block above every outsider. An instance is
\emph{input-prime} if it has no such partition with two or more blocks.
Number a proposer's events from $1$, and put
\[
 F(q)=\{(p,k):1\le k\le q(p)\},\qquad
 \mathcal F(I)=\{F(q):(h,q)\in\Reach(I)\},\qquad
 E(I)=\bigcup\mathcal F(I).
\]
The change from zero-based event indices in~\cite{ishida} is only a
relabeling. The instance is \emph{behaviorally prime} if no partition
$E(I)=E_1\sqcup E_2$ into nonempty sets satisfies
\[
 \mathcal F(I)=\{X_1\cup X_2:X_j\in\mathcal F(I)|_{E_j},\ j=1,2\},
 \qquad
 \mathcal F(I)|_{E_j}=\{X\cap E_j:X\in\mathcal F(I)\}.
\]
It is \emph{doubly prime} if both properties hold. Let
\[
 m_n=\min\{|\Reach(I)|: I\text{ is doubly prime of order }n\},
 \qquad \min\varnothing=+\infty.
\]

\begin{theorem}\label{thm:main}
For every $n\ge2$, $m_n=3\cdot2^{n-1}-1$. A minimizing instance can
simultaneously have total proposal count $2n-1$.
\end{theorem}

Ishida~\cite[Proposition~21.4 and Section~7(A)]{ishida} proves the lower
bound $2^n+n-1$ and asks for the equality above. The rejection-chain
construction in Section~\ref{sec:sharp} already appears in the supplied
research attempts for catalogue entry C78-352;\footnote{See the
\href{https://github.com/mehmetmars7/Erdosproblems-llm-hunter/blob/00ced55ba1223fb48603e1f15028d753997ee515/attacks/open_problems/economics/gpt_6_astra_ultra/C78-352.tex}{GPT~6 Astra attempt, ``A chain of rejections''}, and the
\href{https://github.com/mehmetmars7/Erdosproblems-llm-hunter/blob/00ced55ba1223fb48603e1f15028d753997ee515/attacks/open_problems/economics/opus_5.5_high/C78-352.tex}{Opus~5.5 attempt, ``The upper bound for every $n$''}, pinned to repository commit \texttt{00ced55ba1223fb48603e1f15028d753997ee515}.
Both leave the general lower bound unresolved.}
we include its proof to
make attainability explicit. The contribution here is the universal
matching lower bound. In particular, we do not use the experimentally
checked converse relating proposal overlap to behavioral factors
discussed in~\cite[Sections~2 and~4]{ishida}.

\section{Counting ideals by component merges}

For a finite directed acyclic graph $D$, an \emph{ideal} is a vertex set
closed under predecessors. Let $\Id(D)$ be its set of ideals. Connectivity
below refers to the underlying undirected graph.

\begin{lemma}\label{lem:dag}
If $D$ is connected, has $n\ge1$ sources, and every nonsource has at most
two incoming neighbors, then
\[
                         |\Id(D)|\ge3\cdot2^{n-1}-1.
\]
\end{lemma}
\begin{proof}
Choose a topological order placing all sources first. After these $n$
vertices have been inserted, the processed graph has $n$ singleton
components. Every later vertex attaches to one or two existing
components. Exactly $n-1$ vertices merge two distinct components, since
the final graph is connected. Every processed component contains a
source.

For a merging vertex $v$, let $s(v)$ be the number of sources in the two
components being merged. Every ancestor of $v$ lies in those components.
Let $A(v)$ consist of $v$ and all its ancestors. For every subset $T$ of
the $n-s(v)$ sources outside those components, $A(v)\cup T$ is an ideal.
These $2^{n-s(v)}$ ideals are distinct, since their intersection with the
outside sources recovers $T$. Their latest nonsource in the fixed
topological order is $v$. Thus the families for different merging
vertices are disjoint, and none consists solely of sources. All $2^n$
subsets of the sources are themselves ideals.

Record each component merger as a binary node, with the original source
components as leaves. Vertices that do not merge components leave this
record unchanged. The result is a full binary tree $T$ with $n$ leaves;
$s(v)$ is the number of descendant leaves of its internal node $v$.
Define
\[
             C(T)=\sum_{v\text{ internal}}2^{n-s(v)}.
\]
We claim $C(T)\ge2^{n-1}-1$. For a single leaf both sides are zero.
If the root has subtrees with $a,b\ge1$ leaves, where $a+b=n$, induction
gives
\begin{align*}
 C(T)&=2^bC(T_a)+2^aC(T_b)+1\\
     &\ge2^n-2^a-2^b+1
      \ge2^{n-1}-1.
\end{align*}
The last inequality follows from $(2^a-2)(2^b-2)\ge0$, equivalently
$2^a+2^b\le2^{n-1}+2$. Adding the source-only ideals proves the lemma.
\end{proof}

\section{Rejection witnesses give reachable ideals}

We first record the execution facts needed to apply
Lemma~\ref{lem:dag}. A receiver always holds its favorite among its past
proposers. Thus $q$ determines $h$, and distinct event sets determine
distinct states. Every run has at most $n^2$ steps. Two enabled steps by
different proposers commute and leave each other enabled: another
proposer's action cannot make an unheld proposer held. If they contact
the same receiver, its holder after both steps is the best of the old
holder and the two newcomers.

Consequently the terminal state is unique. For completeness, induct on
$n^2-\sum_p q(p)$: two distinct first steps have a common successor after
commuting; the inductive uniqueness of terminal descendants then
identifies the terminal states of both branches. Every prefix extends to
this terminal state. Denote its counts by $q^*$ and its events by $E$.
Every proposer is held there. Otherwise an unheld proposer must have
exhausted its list, so all $n$ receivers have received proposals and hold
$n$ distinct proposers, contradicting the existence of an unheld one.
In particular, all $n$ first events occur.

Let $H$ be the graph on proposers joining those who contact a common
receiver among the events in $E$. If $H$ were disconnected, events in
different components would address disjoint sets of receivers. Deleting
the other components' steps from a legal prefix preserves legality;
conversely, any projected prefixes can be interleaved because their
receiver updates are disjoint. Every prefix uses only $E$, since it
extends to the unique terminal state. Hence the full event family is
exactly the product of its component projections, each nonempty.
Behavioral primality therefore implies that $H$ is connected.

\begin{lemma}\label{lem:witness}
For a behaviorally prime order-$n$ instance, there is a connected
directed acyclic graph $D$ on $E$ with exactly $n$ sources and exactly two
incoming neighbors at each nonsource, such that every ideal of $D$ is a
distinct reachable state.
\end{lemma}
\begin{proof}
Fix a complete legal execution and write $e(p,k)$ for its $k$th proposal
by $p$. For each $k>1$, let $r$ be the receiver contacted by
$e(p,k-1)$. Immediately before $e(p,k)$, proposer $p$ is unheld and $r$
holds some $w$ it prefers to $p$. Choose $w$'s earlier proposal $e(w,j)$
to $r$ as a witness, and add the edges
\[
       e(p,k-1)\longrightarrow e(p,k),\qquad
       e(w,j)\longrightarrow e(p,k).
\]
Both edges point forward in the fixed execution. Their tails are
distinct because $w\ne p$. The sources are exactly the first events.

Let $X$ be an ideal and execute its events in topological order.
Inductively suppose the preceding events form a legal prefix. The
own-proposer edges ensure that, when $e(p,k)$ is next, $p$ has made
exactly its first $k-1$ proposals. For $k=1$ it is unheld. For $k>1$, the
witness has already proposed to $p$'s previous receiver, whose holder
only improves; that receiver therefore cannot hold $p$. Neither can an
earlier receiver: $p$ was unheld before its following proposal, and
another proposer's action cannot make it held again. Thus $p$ is unheld
and its next proposal is legal. This proves reachability of $X$.
Distinct ideals give distinct proposal counts and hence distinct states.

It remains to prove connectivity. Contract each proposer's own event
chain. Fix a receiver $r$. Every proposer $p$ who contacted $r$ except
its final holder has a subsequent own proposal: otherwise $p$ would be
unheld at termination. Its chosen witness at $r$ is some proposer $w$
strictly preferred to $p$, and the contracted graph contains $p w$.
Iterating such strictly improving witness links among the finite set of
proposers to $r$ reaches its best proposer, the final holder. Therefore
all proposers who contacted $r$ are connected in the contraction. This
connects the endpoints of every edge of $H$. Since $H$ is connected, so
are the contraction and $D$.
\end{proof}

The witness graph need not describe every legal prefix: choosing one
witness selects sufficient prerequisites. Its ideals form a subfamily
of the reachable event family, which is all the lower bound requires.
Lemmas~\ref{lem:dag} and~\ref{lem:witness} prove
\[
       |\Reach(I)|\ge3\cdot2^{n-1}-1
\]
for every behaviorally prime instance, and in particular every doubly
prime instance. The same merger argument gives $|E|\ge2n-1$: besides
the $n$ sources there are $n-1$ distinct merging vertices.

\section{Attainment and double primality}\label{sec:sharp}

For $n\ge2$, take proposers $p_0,p_1,\ldots,p_{n-1}$ and receivers
$r_1,\ldots,r_n$. Let $p_0$ rank $r_1,\ldots,r_n$ in that order.
For $1\le i<n$, let $p_i$ and $r_i$ rank each other first. Complete
every unspecified list arbitrarily to a strict complete order.

Each $p_i$ with $i>0$ makes one proposal and is held permanently at
$r_i$. Only $p_0$ makes later proposals. At $q(p_0)=k-1$, its $k$th
proposal, for $k\ge2$, requires rejection from $r_{k-1}$ by $p_{k-1}$.
It follows that the reachable pointers are exactly
\[
 \begin{array}{ll}
 q(p_0)=0:&q(p_i)\in\{0,1\}\quad(i>0);\\
 q(p_0)=k,\ 1\le k\le n:&q(p_i)=1\quad(1\le i<k),\qquad
 q(p_i)\in\{0,1\}\quad(k\le i<n).
 \end{array}
\]
For sufficiency, execute the selected first proposals of the $p_i$,
then the first $k$ proposals of $p_0$, stopping immediately after the
last one, even if rejected. Consequently
\[
 |\Reach(I_n)|=2^{n-1}+\sum_{k=1}^{n}2^{n-k}
             =3\cdot2^{n-1}-1.
\]

To check behavioral primality directly, name $p_0$'s events
$a_1,\ldots,a_n$ and $p_i$'s event $b_i$. The displayed pointers say
that the event family is precisely the ideals of
\[
          a_1<a_2<\cdots<a_n,\qquad
          b_i<a_{i+1}\quad(1\le i<n).
\]
A mandatory predecessor $e$ of an event $f$ cannot lie in a different
factor of an event-family product: combine a projection containing $f$
with the empty projection in $e$'s factor to contradict that
requirement. All events occur, and the displayed predecessor edges
connect them all. Thus no nontrivial product exists.

For input primality, suppose the block containing $p_0$ has size $k<n$.
Its receivers must be $r_1,\ldots,r_k$, the first $k$ choices of $p_0$.
Since $r_i$ ranks $p_i$ first, that block must contain each
$p_1,\ldots,p_k$ as well as $p_0$, contradicting its size $k$.
Therefore $I_n$ is doubly prime. This proves attainability and completes
Theorem~\ref{thm:main}. Its terminal counts are $q(p_0)=n$ and
$q(p_i)=1$ for $i>0$, giving $2n-1$ proposals. The same instance thus
attains the proposal lower bound of~\cite[Proposition~21.4]{ishida}
and the state minimum, answering the additional joint-attainment
question in~\cite[Section~7(A)]{ishida}.


\section*{Verification, attribution, and scope}
This submission claims a complete hand proof for the catalogue statement,
including all $n\ge2$ and joint attainment with $2n-1$ proposals.
The mathematical proof and Lean companion were
prepared with GPT 6 Astra (configured model \texttt{gpt-6-astra}, reasoning
effort \texttt{xhigh}) using parallel proof construction and independent
adversarial checking. The exact serving revision is not exposed.

Earlier catalogue attempts already proved the matching rejection-chain
upper construction. We preserve and credit the
\href{https://github.com/mehmetmars7/Erdosproblems-llm-hunter/blob/00ced55ba1223fb48603e1f15028d753997ee515/attacks/open_problems/economics/gpt_6_astra_ultra/C78-352.tex}{earlier GPT attempt}
and
\href{https://github.com/mehmetmars7/Erdosproblems-llm-hunter/blob/00ced55ba1223fb48603e1f15028d753997ee515/attacks/open_problems/economics/opus_5.5_high/C78-352.tex}{earlier Opus attempt}.
The new contribution claimed here is the general lower bound, completing
the equality; the upper construction is included for a self-contained proof.

The Lean companion formalizes finite combinatorial counting statements
about binary merger trees and an abstract event-system bound. It does not
formalize the complete deferred-acceptance dynamics, the witness-graph
construction and connectivity bridge, or the complete catalogue theorem.
These bridges are proved in the manuscript. Successful Lean checking of
the companion alone would not establish the full claim.

The pinned environment is Lean 4.35.0-rc3 with Mathlib commit
\texttt{b84a70d6a5ed793cc46184160f4d4188d8c825fb}.

The source package is identified by the following commit-pinned references:
\href{https://github.com/jbaelaw/prime-da-state-counts/blob/eec79a6a79c77683811142c06c7c5cc71ad48bf6/paper/Proof.pdf}{four-page proof PDF},
\href{https://github.com/jbaelaw/prime-da-state-counts/blob/eec79a6a79c77683811142c06c7c5cc71ad48bf6/paper/Proof.tex}{complete manuscript source},
\href{https://github.com/jbaelaw/prime-da-state-counts/tree/eec79a6a79c77683811142c06c7c5cc71ad48bf6/provenance}{complete original generated proof and independent audit outputs},
and \href{https://github.com/jbaelaw/prime-da-state-counts/tree/eec79a6a79c77683811142c06c7c5cc71ad48bf6/lean}{Lean companion}.
The principal checked statements are \texttt{ideals\_card\_lower} in
\href{https://github.com/jbaelaw/prime-da-state-counts/blob/eec79a6a79c77683811142c06c7c5cc71ad48bf6/lean/BoundIdeal.lean}{BoundIdeal.lean}, and
\texttt{weight\_lower}, \texttt{card\_witness\_lower},
\texttt{state\_count\_of\_injection}, and \texttt{card\_witness\_comb} in
\href{https://github.com/jbaelaw/prime-da-state-counts/blob/eec79a6a79c77683811142c06c7c5cc71ad48bf6/lean/BoundTree.lean}{BoundTree.lean}.

The remote \texttt{lake build} and \texttt{lake env lean Verify.lean}
checks passed. The audited declarations depend only on
\texttt{propext}, \texttt{Classical.choice}, and \texttt{Quot.sound};
no \texttt{sorryAx}, admitted proofs, or additional axioms occur.
The compiler emitted two module-system compatibility notices and two
unused section-variable notices, with no proof errors.
\href{https://github.com/jbaelaw/prime-da-state-counts/blob/eec79a6a79c77683811142c06c7c5cc71ad48bf6/verification/LeanVerification.json}{The exact formalization scope and checking record}
and \href{https://github.com/jbaelaw/prime-da-state-counts/tree/eec79a6a79c77683811142c06c7c5cc71ad48bf6/lean/verification}{logs and source hashes}
are retained. The full deferred-acceptance theorem remains a hand-proof
claim, with the formalization boundary stated above.

The \href{https://github.com/jbaelaw/prime-da-state-counts/blob/eec79a6a79c77683811142c06c7c5cc71ad48bf6/verification/verify.py}{independent finite harness}
passed its default run: 683 directed acyclic graphs, including 396
connected graphs; the explicit family for orders two through eight;
all 16 strict order-two profiles; and 3,464 replays of witness ideals.
\href{https://github.com/jbaelaw/prime-da-state-counts/blob/eec79a6a79c77683811142c06c7c5cc71ad48bf6/verification/verification.json}{Exact results}
and \href{https://github.com/jbaelaw/prime-da-state-counts/blob/eec79a6a79c77683811142c06c7c5cc71ad48bf6/paper/CitationAudit.json}{citation checks}
are available. These finite checks supplement the arbitrary-order
hand proof. The PDF compiled successfully and all four pages were
visually inspected. No statement here implies that the catalogue
maintainers have verified or endorsed the proof.

\section*{Completion Estimate}
100\%

This is the estimate for the hand-proof claim; the Lean companion is partial in the
precise sense stated above. This is an author-submitted LLM proof claim
awaiting independent expert review, not a claim of verification by the
catalogue maintainers.

\begin{thebibliography}{9}
\bibitem{ishida}
Y.~Ishida.
\emph{Prime factorisation of stable-matching instances: uniqueness,
simultaneous products, and an exact census}.
arXiv:2610.05617v1, 2026.
\url{https://arxiv.org/html/2610.05617v1}.
\end{thebibliography}
\end{document}
